%% 04_evaluacion.tex — Resultados de la barrida de ablación.
%%
%% Contenido numérico: testbed/experiment_results.csv (graduado por el
%% oráculo) + testbed/analysis/{summary_by_run.csv,stats.json,early_recall.json}
%% + figuras de paper/figures/, todo producido por tools/analyze_results.py.
%% Las cifras en prosa provienen de sections/_datos.tex.

\section{Evaluación y resultados}
\label{sec:resultados}

La barrida produjo \NumCorridas{}~ejecuciones calificadas: la rejilla completa
de los cuatro presupuestos $\{10, 20, 50, \infty\}$ con las
\NumCondiciones{}~configuraciones (\code{baseline}, \code{no-interleave},
\code{no-priority}, \code{no-runtime-confirm} y \code{no-adaptive-sorting}).
Todos los valores proceden de \code{tools/analyze\_results.py}.

Adelantamos la lectura global: la barrida de ablación es en su mayor parte
\emph{nula} (\cref{sec:res-stats}), y lo que sí tiene valor práctico no es qué
etapa se apaga sino \emph{cuántas peticiones se gastan}. Organizamos la sección
en torno a esa \emph{economía de recursos}.

\subsection{La economía del presupuesto: dónde está el punto óptimo}
\label{sec:res-eficiencia}

La \cref{fig:deteccion-vs-presupuesto} y la \cref{tab:auc} muestran el patrón
central. Con \num{10}~peticiones por parámetro el \code{baseline} ya recupera
el \ValorRecallBDiez{} de las instancias vulnerables; con \num{20} alcanza su
valor de \emph{saturación}, \ValorRecallTecho{}, y a partir de ahí
\textbf{ninguna petición adicional detecta una sola vulnerabilidad nueva}:
los presupuestos \num{50} y $\infty$ producen exactamente los mismos
\ValorTPtecho{}~verdaderos positivos que el presupuesto \num{20}.

Lo que sí cambia al ampliar el presupuesto es la precisión, y para peor: los
falsos positivos crecen de \ValorFPBVeinte{} (presupuesto \num{20}) a
\ValorFPBCincuenta{} (presupuesto \num{50}), de modo que la precisión cae
$\DeltaPrecisionSaturacion{}$ (\ValorPrecisionBVeinte{} $\to$
\ValorPrecisionBCincuenta{}) y el $F_1$ $\DeltaFOneSaturacion{}$
(\ValorFOneBVeinte{} $\to$ \ValorFOneBCincuenta{}). El $F_1$ describe así una
curva con máximo claro en \num{20}~peticiones por parámetro
(\ValorFOneBVeinte{}), limitada por exhaustividad con \num{10}
(\ValorFOneBDiez{}) y por precisión con \num{50} o más. Sobre objetivos
estáticos, más \emph{payloads} sobre un punto de inyección ya explorado no
compran cobertura, solo acumulan coincidencias espurias: el presupuesto no
conviene maximizarlo, sino \emph{ajustarlo}.

\begin{figure}[t]
  \centering
  \IfFileExists{figures/fig1_detection_vs_budget.pdf}{%
    \includegraphics[width=\columnwidth]{fig1_detection_vs_budget.pdf}%
  }{%
    \fbox{\parbox[c][3.2cm][c]{0.9\columnwidth}{\centering\footnotesize
      \textsc{Marcador de posición} ---
      \code{figures/fig1\_detection\_vs\_budget.pdf}}}%
  }
  \caption{Exhaustividad (instancias vulnerables únicas detectadas $/$ total)
  en función del presupuesto de peticiones \code{-{}-max-payloads}
  ($\infty$ = sin límite), una curva por configuración del \emph{pipeline}.
  Todas las curvas saturan en \ValorRecallTecho{} a partir de
  \num{20}~peticiones; \code{baseline}, \code{no-interleave},
  \code{no-runtime-confirm} y \code{no-adaptive-sorting} se solapan.}
  \label{fig:deteccion-vs-presupuesto}
\end{figure}

\begin{figure}[t]
  \centering
  \IfFileExists{figures/fig2_ablation_impact.pdf}{%
    \includegraphics[width=\columnwidth]{fig2_ablation_impact.pdf}%
  }{%
    \fbox{\parbox[c][3.2cm][c]{0.9\columnwidth}{\centering\footnotesize
      \textsc{Marcador de posición} ---
      \code{figures/fig2\_ablation\_impact.pdf}}}%
  }
  \caption{Caída de exhaustividad ($\text{baseline} - \text{ablación}$) al
  desactivar cada optimización, por presupuesto. La $\Delta$ FPR es
  idénticamente nula en todas las celdas y se omite de la gráfica. El único
  efecto apreciable es \code{no-priority} en el presupuesto de \num{20}.}
  \label{fig:impacto-ablacion}
\end{figure}

\subsection{El balance de precisión bajo escasez: la priorización por confianza}
\label{sec:res-auc}

El régimen interesante es el \emph{escaso}, y ahí una etapa de planificación
demuestra su valor. A presupuesto \num{20}, el \code{baseline} alcanza un AUC
de \ValorAUCBaseline{} (\cref{eq:auc}), con precisión
\ValorPrecisionBaseline{}, exhaustividad \ValorRecallBaseline{} y $F_1$
\ValorFOneBaseline{}. Desactivar el entrelazado por contexto, la confirmación
en tiempo de ejecución o la ordenación adaptativa no cambia \emph{ninguna} de
esas métricas (idénticas al \code{baseline} hasta la tercera cifra; AUC
\ValorAUCNoInterleave{}, \ValorAUCNoRuntimeConfirm{}, \ValorAUCNoAdaptive{}).

Desactivar la \textbf{priorización por confianza} sí las degrada: el AUC baja a
\ValorAUCNoPriority{}, la exhaustividad cae \DeltaRecallNoPriority{} y el $F_1$
pierde \DeltaFOneNoPriority{} (\ValorFOneBaseline{} $\to$
\ValorFOneNoPriorityBVeinte{}), es decir
\NoPriorityVulnPerdidas{}~vulnerabilidades no detectadas y
\NoPriorityFPnuevos{}~falsos positivos adicionales: sin priorización, el
presupuesto se consume en \emph{payloads} de bajo rendimiento antes de llegar
a los que discriminan. El efecto es exclusivo de la escasez: a partir de
\num{50} las \NumCondiciones{}~configuraciones convergen (precisión dentro de
\num{0.3}~p.p., \cref{tab:auc}). La FPR es \ValorFPRBaseline{} (\num{2} de
\NumTrampas{}~trampas) en toda configuración y presupuesto $\ge \num{20}$: la
priorización cambia \emph{cuáles} \emph{payloads} se prueban, no la tasa a la
que las trampas activan una alarma.

\begin{table}[t]
  \centering
  \caption{Métricas por ejecución y AUC de la curva de detección
  (\cref{eq:auc}). El $F_1$ máximo del \code{baseline} está en el presupuesto
  \num{20} (en negrita). Fuente: \code{summary\_by\_run.csv}.}
  \label{tab:auc}
  \footnotesize
  \setlength{\tabcolsep}{4pt}
  \begin{tabular}{ll S[table-format=1.3] S[table-format=1.3] S[table-format=1.3] S[table-format=1.3]}
    \toprule
    Presup. & Config. & {Prec.} & {Rec.} & {$F_1$} & {AUC} \\
    \midrule
    \multirow{5}{*}{10}
      & baseline            & 0.626 & 0.613 & 0.620 & 0.329 \\
      & no-interleave       & 0.623 & 0.613 & 0.618 & 0.329 \\
      & no-priority         & 0.623 & 0.613 & 0.618 & 0.330 \\
      & no-runtime-confirm  & 0.626 & 0.613 & 0.620 & 0.329 \\
      & no-adaptive-sorting & 0.626 & 0.613 & 0.620 & 0.330 \\
    \midrule
    \multirow{5}{*}{20}
      & baseline            & 0.622 & 0.644 & \bfseries 0.633 & 0.338 \\
      & no-interleave       & 0.622 & 0.644 & 0.633 & 0.338 \\
      & no-priority         & 0.564 & 0.613 & 0.588 & 0.332 \\
      & no-runtime-confirm  & 0.619 & 0.644 & 0.631 & 0.338 \\
      & no-adaptive-sorting & 0.619 & 0.644 & 0.631 & 0.338 \\
    \midrule
    50 & \emph{todas} ($\pm$)   & 0.572 & 0.644 & 0.606 & 0.340 \\
    $\infty$ & \emph{todas} ($\pm$) & 0.573 & 0.644 & 0.607 & 0.340 \\
    \bottomrule
  \end{tabular}
  \\[2pt]
  \footnotesize\emph{Presupuestos $\ge \num{50}$}: las
  \NumCondiciones{}~configuraciones coinciden hasta $\pm\num{0.002}$ en cada
  métrica; se muestra el valor medio.
\end{table}

\subsection{El techo de recuperación: una brecha de cobertura, no de planificación}
\label{sec:res-techo}

El hecho más rotundo de la \cref{tab:auc} es un negativo: los verdaderos
positivos se estancan estrictamente en \ValorTPtecho{}~de las
\NumVulnerables{}~instancias vulnerables, y los falsos negativos nunca bajan de
\ValorFNtecho{}, \emph{con independencia del presupuesto y de la
configuración}. Ninguna combinación de presupuesto ampliado, entrelazado,
priorización, confirmación u ordenación adaptativa recupera una sola de esas
\ValorFNtecho{}~instancias.

La limitación hay que atribuirla bien. No es un fallo del \emph{scheduling}:
el planificador solo decide el \emph{orden} de emisión de un corpus dado, y a
presupuesto $\infty$ lo emite entero. Si aun así \ValorFNtecho{}~instancias
quedan sin detectar, la causa está en el \textbf{corpus mismo}: para esas
tuplas \code{(url, parámetro, tipo)} ninguna carga dispara la condición
vulnerable bajo la codificación o el contexto de salida de ese \emph{endpoint}.
Es una brecha de \emph{cobertura de vectores}, y su solución es ampliar la
biblioteca de \emph{payloads}, no refinar una ordenación que sobre esas
instancias no tiene nada que ordenar (\cref{sec:conclusion}).

\subsection{Significancia estadística}
\label{sec:res-stats}

Sobre los vectores de coste de detección por instancia (bloque agrupado de los
\NumBloquesCompletos{}~presupuestos, $n = \ValorFriedmanN{}$), la prueba de
Friedman rechaza la equivalencia de las cuatro configuraciones de ordenación
($\chi^2 = \ValorFriedmanChiSq{}$, $p = \ValorFriedmanP{}$). Pero el tamaño del
efecto es \emph{insignificante} en las tres comparaciones pareadas
\code{baseline} frente a ablación: $\delta$ de Cliff de
\ValorCliffNoInterleave{} (\code{no-interleave}), \ValorCliffNoPriority{}
(\code{no-priority}) y \ValorCliffNoRuntimeConfirm{}
(\code{no-runtime-confirm}), todos con $\lvert\delta\rvert < \num{0.02}$
(\cref{tab:stats}). El post-hoc de Nemenyi ni siquiera separa
\code{no-runtime-confirm} del \code{baseline}
($p = \ValorNemenyiNoRuntimeConfirm{}$).

Las diferencias que Friedman detecta son de \emph{ordenación fina} de las
peticiones: reales por el gran número de bloques, pero sin traducción práctica
en el conjunto detectado ni en su precisión. El contraste entre la ordenación
adaptativa en vivo (\code{baseline}) y la estática \emph{a priori}
(\code{no-adaptive-sorting}) lo confirma desde el otro lado: la adaptativa
llega a cada vulnerabilidad, en promedio, un \ValorAdaptivePctLento{}
\emph{más tarde} ($\delta$ de Cliff \ValorAdaptiveCliff{}, Wilcoxon
$p = \ValorAdaptiveP{}$), con idéntica exhaustividad final. El único cambio con
consecuencias operativas sigue siendo la pérdida de selectividad de
\code{no-priority} bajo escasez (\cref{sec:res-auc}).

\begin{table}[t]
  \centering
  \caption{Pruebas sobre el coste de detección por instancia, \code{baseline}
  frente a cada ablación (bloque agrupado, $n = \ValorFriedmanN{}$). Wilcoxon
  con corrección de Benjamini--Hochberg; $p_{\text{Nem.}}$ del post-hoc de
  Nemenyi; $\delta$ de Cliff como tamaño del efecto. Fuente:
  \code{stats.json}.}
  \label{tab:stats}
  \footnotesize
  \setlength{\tabcolsep}{4pt}
  \resizebox{\columnwidth}{!}{%
  \begin{tabular}{lcccl}
    \toprule
    \code{baseline} vs & $p_{\text{BH}}$ & $p_{\text{Nem.}}$
      & $\delta$ Cliff & Magnitud \\
    \midrule
    no-interleave      & \ValorWilcoxonPNoInterleave{}
      & \ValorNemenyiNoInterleave{} & \ValorCliffNoInterleave{}
      & \MagnitudCliffNoInterleave{} \\
    no-priority        & \ValorWilcoxonPNoPriority{}
      & \ValorNemenyiNoPriority{} & \ValorCliffNoPriority{}
      & \MagnitudCliffNoPriority{} \\
    no-runtime-confirm & \ValorWilcoxonPNoRuntimeConfirm{}
      & \ValorNemenyiNoRuntimeConfirm{} & \ValorCliffNoRuntimeConfirm{}
      & \MagnitudCliffNoRuntimeConfirm{} \\
    \bottomrule
  \end{tabular}}
  \\[2pt]
  \footnotesize\emph{Friedman global: $\chi^2 = \ValorFriedmanChiSq{}$,
  $p = \ValorFriedmanP{}$.}
\end{table}

El diagrama de diferencias críticas correspondiente (post-hoc de Nemenyi sobre
el rango medio de coste de detección, \code{fig3\_cd\_diagram} en el
repositorio) sitúa a \code{no-priority} en el mejor rango, pero solo porque
adelanta la primera detección probando antes \emph{payloads} triviales de alto
rendimiento aparente; como muestra la \cref{tab:auc}, lo paga con
\NoPriorityVulnPerdidas{}~vulnerabilidades y \NoPriorityFPnuevos{}~falsos
positivos a presupuesto \num{20}. Con todos los tamaños del efecto
insignificantes, el orden de rangos es indicativo, no una ventaja práctica.

\subsection{Síntesis de hallazgos}
\label{sec:res-sintesis}

\begin{enumerate}
  \item \textbf{Punto óptimo de presupuesto.} El $F_1$ del \code{baseline} es
  máximo en \num{20}~peticiones por parámetro (\ValorFOneBVeinte{}); con
  \num{10} lo limita la exhaustividad (\ValorFOneBDiez{}) y con \num{50}$+$ la
  precisión (\ValorFOneBCincuenta{}). Pasar de \num{20} a \num{50}$+$ no aporta
  \emph{ni un} verdadero positivo nuevo pero suma
  \SaturacionFPnuevos{}~falsos positivos: sobre \emph{endpoints} estáticos,
  más \emph{payloads} son ruido, no cobertura.
  \item \textbf{La priorización por confianza es la única etapa con efecto
  práctico, y solo bajo escasez} ($-\DeltaRecallNoPriority{}$ de exhaustividad
  y $-\DeltaPrecisionNoPriority{}$ de precisión a presupuesto \num{20}; nada a
  partir de \num{50}). El entrelazado, la confirmación en tiempo de ejecución
  y la ordenación adaptativa no alteran exhaustividad, precisión ni FPR: su
  efecto sobre el orden es detectable pero insignificante.
  \item \textbf{El techo de \ValorTPtecho{}$/$\NumVulnerables{} es
  estructural.} Las \ValorFNtecho{}~instancias no detectadas lo siguen estando
  a presupuesto $\infty$: brecha de cobertura del corpus, ajena al
  planificador.
\end{enumerate}
